% ============================================================================
%  Subdocumento 5. Contenido. Se usa igual en el documento único y en el
%  PDF propio del subdocumento. Los formularios que le asigna el T-21 se
%  agregan solos al final (datos en formularios/ de esta carpeta).
% ============================================================================

\begin{resumenApertura}
\marcador{Qué resuelve este subdocumento, en cuatro o cinco líneas}
\begin{recibe}
  \item \marcador{Qué recibe el mandante}
\end{recibe}
\end{resumenApertura}

\section{Dominio de información}\label{sec:5-dominio-de-informacion}

\marcador{Entidades y agregados, incluida la evidencia de jornada}

\section{Motor y paradigma de persistencia}\label{sec:5-motor-y-paradigma-de-persistenci}

\marcador{Motor, paradigma, transaccionalidad y posición CAP}

\section{Migración, saneamiento y conciliación}\label{sec:5-migracion-saneamiento-y-concilia}

\marcador{Estrategia de migración de los datos históricos}

\section{Estrategia de desempeño}\label{sec:5-estrategia-de-desempeno}

\marcador{Indexación, particionamiento, caché y optimización de consultas}

\section{Almacenamiento transaccional y analítico}\label{sec:5-almacenamiento-transaccional-y-a}

\marcador{Separación y modelo de explotación}

\section{Calidad, retención, archivado y eliminación}\label{sec:5-calidad-retencion-archivado-y-el}

\marcador{Reglas por dominio de datos}
